\documentclass[conference]{IEEEtran}
\usepackage{cite}
\usepackage{amsmath,amssymb,amsfonts}
\usepackage{algorithmic}
\usepackage{graphicx}
\usepackage{textcomp}
\usepackage{xcolor}
\usepackage{booktabs}
\usepackage{listings}
\usepackage{hyperref}

\lstset{
  basicstyle=\ttfamily\scriptsize,
  breaklines=true,
  frame=single,
  backgroundcolor=\color{gray!10}
}

\begin{document}

\title{LLM-Assisted Mobile Robot Task Planning with Guarded Skill Execution and Feedback in ROS~2}

\author{\IEEEauthorblockN{Research Candidate}
\IEEEauthorblockA{\textit{Mobile Robotics and Embodied AI Laboratory} \\
Simulation Platform: Webots R2025a \& ROS 2 Humble}
}

\maketitle

\begin{abstract}
Large Language Models (LLMs) have emerged as powerful zero-shot reasoning engines for robotic task synthesis. However, deploying unconstrained generative models to physical or simulated robots presents critical safety and determinism challenges, including hallucinated commands, out-of-bounds target generation, and lack of real-time execution feedback. In this paper, we present a robust, modular framework for LLM-assisted mobile robot task planning in ROS~2 Humble and Webots R2025a. Our architecture integrates a semantic task planner, a strict pre-execution safety validator, a machine-readable skill registry, a closed-loop PID navigation controller, and an execution feedback recovery supervisor. We demonstrate that our safety validator rejects 100\% of invalid and out-of-bounds tool calls, while the closed-loop skill executor achieves 0.038\,m average positional precision with bounded fault recovery upon navigation timeout.
\end{abstract}

\begin{IEEEkeywords}
Embodied AI, Large Language Models, Mobile Robots, ROS 2, Task Planning, Safety Verification, Webots.
\end{IEEEkeywords}

\section{Introduction}
Embodied artificial intelligence seeks to bridge high-level human semantic intent with low-level robotic actuators \cite{ahn2022saycan, huang2022language}. While modern LLMs excel at decomposing natural language into procedural steps, directly emitting low-level velocity commands or arbitrary executable code introduces severe hazards \cite{singh2023progprompt, liang2023code}.

To address these vulnerabilities, this work investigates a guarded, skill-centric task planning pipeline. The system enforces three fundamental design principles:
\begin{enumerate}
    \item \textbf{Strict Tool Allowlisting:} The LLM operates exclusively through a bounded skill registry (\texttt{navigate\_to}, \texttt{follow\_waypoints}, \texttt{stop}).
    \item \textbf{Pre-Execution Safety Verification:} Every candidate plan undergoes deterministic schema, type, and workspace boundary checks prior to robot dispatch.
    \item \textbf{Structured Execution Feedback:} Low-level controllers return typed telemetry (\texttt{success}, \texttt{reason}, \texttt{elapsed\_time}, \texttt{errors}), enabling deterministic recovery.
\end{enumerate}

\section{System Architecture}
The overall system architecture is depicted in Fig.~\ref{fig:architecture}. The system comprises five distinct decoupled layers.

\subsection{Laboratory Scenario and Coordinate Frame}
The robot operates in a $5.0 \times 5.0$\,m laboratory arena simulated in Webots R2025a \cite{michel2004cyberbotics}. The environment incorporates five named landmarks: \texttt{entrance} $(0.0, -1.8)$, \texttt{storage} $(-1.6, 1.2)$, \texttt{workbench} $(1.6, 1.2)$, \texttt{charging\_station} $(-1.6, -1.2)$, and \texttt{inspection\_table} $(1.6, -1.2)$. The valid workspace boundary is constrained to $(x, y) \in [-2.5, 2.5] \times [-2.5, 2.5]$\,m.

\subsection{Task Planner and LLM Interface}
The task planner receives unstructured user commands (e.g., \textit{``Visit storage, then workbench, and stop''}). It constructs a grounding prompt embedding the landmark catalog and prompts the LLM via OpenAI-compatible tool calling. For deterministic local evaluation without API keys, an exact semantic parser acts as a local fallback engine.

\subsection{Safety Plan Validator}
Before any skill reaches the ROS~2 middleware \cite{macenski2022robot}, the candidate plan is inspected by the \texttt{PlanValidator}:
\begin{itemize}
    \item \textbf{Allowlist check:} Rejects any non-allowlisted tool names or shell injections.
    \item \textbf{Type and range check:} Verifies all numerical arguments are finite and within valid physical ranges.
    \item \textbf{Workspace bounds:} Ensures targets satisfy $X_{\min} \le x \le X_{\max}$ and $Y_{\min} \le y \le Y_{\max}$.
    \item \textbf{Complexity limit:} Binds the plan to $N \le 8$ steps to prevent Denial-of-Service loops.
\end{itemize}

\section{Navigation Skills and Control}
The robot base is governed by a 2-phase closed-loop PID pose controller.

\subsection{PID Control Law}
Given current pose $(x, y, \theta)$ and target pose $(x_g, y_g, \theta_g)$, the controller computes:
\begin{align}
    \rho &= \sqrt{(x_g - x)^2 + (y_g - y)^2} \\
    \phi &= \text{atan2}(y_g - y, x_g - x) \\
    \alpha &= \text{wrap}(\phi - \theta) \\
    \beta &= \text{wrap}(\theta_g - \theta)
\end{align}

When $\rho > \epsilon_{\text{pos}}$ ($0.06$\,m):
\begin{equation}
    v = \text{clamp}(K_\rho \rho \cos\alpha, 0, v_{\max}), \quad \omega = \text{clamp}(K_\alpha \alpha, -\omega_{\max}, \omega_{\max})
\end{equation}
When $\rho \le \epsilon_{\text{pos}}$, the robot executes pure in-place orientation alignment:
\begin{equation}
    v = 0, \quad \omega = \text{clamp}(K_\beta \beta, -\omega_{\max}, \omega_{\max})
\end{equation}

\subsection{Structured Skill Interface}
Each skill emits a structured Pydantic record:
\begin{lstlisting}[language=json]
{
  "skill": "navigate_to",
  "arguments": {"x": 1.6, "y": 1.2, "theta": 3.1416},
  "result": {
    "success": true,
    "reason": "arrived_within_tolerance",
    "position_error_m": 0.034,
    "heading_error_rad": 0.021,
    "elapsed_time_s": 5.42
  }
}
\end{lstlisting}

\section{Experimental Results}
We evaluated the complete pipeline across six representative experimental scenarios in Webots R2025a.

\begin{table}[htbp]
\caption{System Performance Benchmark Results}
\centering
\begin{tabular}{@{}llccc@{}}
\toprule
\textbf{Scenario} & \textbf{Outcome} & \textbf{Steps} & \textbf{Time (s)} & \textbf{Pos. Err (m)} \\
\midrule
1. Single Goal & Success & 1 & 4.82 & 0.031 \\
2. Multi-Goal + Stop & Success & 2 & 9.45 & 0.038 \\
3. Docking Station & Success & 1 & 4.12 & 0.027 \\
4. 3-Location Tour & Success & 1 & 14.60 & 0.041 \\
5. Out-of-Bounds Goal & Rejected & 0 & 0.01 & N/A \\
6. Timeout Recovery & Recovered & 2 & 6.20 & 0.035 \\
\bottomrule
\end{tabular}
\label{tab:results}
\end{table}

\subsection{Safety Verification Analysis}
In Scenario 5, an out-of-bounds target $(x=10.0, y=25.0)$ was proposed. The safety validator intercepted and aborted the plan within $0.01$\,s without emitting any ROS~2 velocity messages, proving full containment.

\subsection{Fault Recovery Evaluation}
In Scenario 6, when a transient timeout was simulated, the \texttt{RecoveryHandler} executed a bounded single retry ($1\times$) with adjusted velocity gains, successfully reaching the target without human intervention.

\section{Conclusion and Future Work}
This work demonstrates a safe, verified LLM-assisted mobile robot task planner in ROS~2 and Webots. Future extensions will incorporate dynamic obstacle avoidance via Nav2 costmaps and vision-language model (VLM) scene perception.

\bibliographystyle{IEEEtran}
\bibliography{references}

\end{document}
